\begin{statementbox}
\textbf{\textcolor{CStatement}{条件与记号}}
\begin{description}[style=nextline, leftmargin=2.8em, labelsep=0.5em, itemsep=0.35em]
    \item[\textbf{条件 1（椭圆）：}] $E:\dfrac{x^2}{a^2}+\dfrac{y^2}{b^2}=1$，$a>b>0$；设 $c=\sqrt{a^2-b^2}$，焦点为 $F_1(-c,0)$、$F_2(c,0)$。
    \item[\textbf{条件 2（直线）：}] $l:y=kx+m$，其中 $k,m\in\mathbb R$，即直线斜率存在。
\end{description}

\medskip
\textbf{\textcolor{CStatement}{核心结论（充要条件）}}
\[
\boxed{l\text{ 与 }E\text{ 相切}\iff m^2=a^2k^2+b^2.}
\]
\textbf{固定斜率 $k$ 时，两条切线：}
\[
\boxed{y=kx\pm\sqrt{a^2k^2+b^2}.}
\]
对应切点（当 $m^2=a^2k^2+b^2$ 时）为
\[
\boxed{P\left(-\frac{a^2k}{m},\frac{b^2}{m}\right).}
\]
由于 $m^2=a^2k^2+b^2>0$，切点公式的分母 $m$ 不会为零。

\medskip
\textbf{\textcolor{CStatement}{位置关系（顺带一判三）}}
令 $D=a^2k^2+b^2$，则
\[
\begin{array}{ccl}
 m^2<D &\Longleftrightarrow& \text{相交于两个不同点（割线）},\\
 m^2=D &\Longleftrightarrow& \text{相切于一个点},\\
 m^2>D &\Longleftrightarrow& \text{无公共点（相离）}.
\end{array}
\]

\medskip
\textbf{\textcolor{CStatement}{焦点距离性质（重要限定）}}
若 $l$ 是切线，$F_1,F_2$ 到 $l$ 的距离分别为 $d_1,d_2$，则
\[
\boxed{d_1d_2=b^2.}
\]
\textbf{注意：}单有 $d_1d_2=b^2$ \emph{不能}推出相切；必须再要求两焦点位于 $l$ 的同侧，才与相切等价。

\medskip
\textbf{\textcolor{CStatement}{特殊情形}}
斜率不存在的两条切线为 $\boxed{x=\pm a}$；当 $k=0$ 时，上述公式给出 $y=\pm b$。
\end{statementbox}
